\section*{Resumen.}

El presente proyecto de ingeniería de software expone el diseño, desarrollo e implementación de \textit{Vitalis PharmaSys}, una plataforma web integral orientada a la gestión y control operativo de establecimientos farmacéuticos independientes y comunitarios. El objetivo fundamental del sistema es sistematizar la administración de inventarios, optimizar la dispensación de medicamentos y garantizar la trazabilidad sanitaria de productos sujetos a prescripción médica, mitigando las pérdidas económicas por caducidad y minimizando los errores operativos derivados de registros manuales. La solución tecnológica se fundamenta en el patrón arquitectónico Modelo-Vista-Controlador (MVC), desarrollada mediante el framework Laravel sobre el lenguaje PHP en su versión 8.4, articulada con un sistema gestor de bases de datos relacionales MariaDB bajo transacciones ACID, y maquetada con Tailwind CSS bajo los lineamientos del sistema de diseño \textit{Clinical Precision UI}. Entre sus módulos funcionales destacan la gestión de existencias mediante la estrategia logística FEFO (\textit{First Expired, First Out}), la emisión automatizada de alertas preventivas por caducidad y desabasto, el control riguroso de recetas médicas para antibióticos y psicotrópicos, un punto de venta ágil y una bitácora inmutable de auditoría operacional. Los resultados obtenidos validan la eficacia del sistema al reducir drásticamente el tiempo de atención en mostrador, erradicar discrepancias en el inventario físico y asegurar el cumplimiento riguroso de las regulaciones sanitarias vigentes, proporcionando a los administradores una herramienta confiable para la toma de decisiones estratégicas.
